\documentclass[12pt]{article}
% Préambule commun à tous les documents extraits de « La Revue CAPES »
\usepackage[T1]{fontenc}
\usepackage[utf8]{inputenc}
\usepackage{lmodern}
\usepackage{amsmath,amssymb,amsthm,amsfonts,mathrsfs}
\usepackage{setspace}
\usepackage{listings}
\usepackage{pgfplots}
\pgfplotsset{compat=1.18}
\usepackage{longtable}
\usepackage{adjustbox}
\usepackage{lettrine}
\usepackage{tkz-tab}
\usepackage{changepage}
\usepackage{pdflscape}
\usepackage{graphicx}
\usepackage{tikz}
\usepackage{xcolor}
\usepackage[a4paper,top=2cm,bottom=2.5cm,left=2.5cm,right=2cm]{geometry}
\usepackage{fancyhdr}
\usepackage{draftwatermark}
\SetWatermarkText{La RevueCapes}
\SetWatermarkLightness{0.9}
\SetWatermarkScale{2}
\usepackage{hyperref}
\hypersetup{colorlinks=true,linkcolor=blue!50!black,urlcolor=blue!50!black}

\definecolor{revue}{RGB}{20,60,120}

\pagestyle{fancy}
\fancyhf{}
\renewcommand{\headrulewidth}{0.4pt}
\setlength{\headheight}{15pt}

% \entete{Matière}{Titre}{Type de document}
\newcommand{\entete}[3]{%
  \fancyhead[L]{\small\textsc{La Revue CAPES} -- #1}%
  \fancyhead[R]{\small #2}%
  \fancyfoot[C]{\thepage}%
  \thispagestyle{plain}%
  \begin{center}
    {\color{revue}\rule{\textwidth}{1.2pt}}\\[4pt]
    {\small\textsc{Concours directs CAP-CEG / CAPES -- Burkina Faso}}\\[6pt]
    {\LARGE\bfseries\color{revue} #2}\\[6pt]
    {\large\itshape #1 -- #3}\\[2pt]
    {\color{revue}\rule{\textwidth}{1.2pt}}
  \end{center}
  \vspace{0.5cm}%
}

% Encadré pour les remarques ajoutées dans les corrigés
\newcommand{\remarque}[1]{\par\medskip\noindent\fcolorbox{revue}{revue!6}{\parbox{\dimexpr\linewidth-2\fboxsep-2\fboxrule}{\textbf{\color{revue}Remarque.} #1}}\par\medskip}


\begin{document}
\entete{Mathématiques}{Corrigé détaillé du sujet type n°11}{Correction}
\begin{spacing}{1.5}

\textbf{Exercice 1}

\emph{Rappel.} Si $f$ est continue sur $[a,+\infty[$, l'intégrale $\int_a^{+\infty}f(t)\,dt$ converge si et seulement si $x\mapsto\int_a^xf(t)\,dt$ admet une limite finie en $+\infty$ ; cette limite est alors la valeur de l'intégrale.

\textbf{1. Calcul d'intégrales généralisées.}

\textbf{a.} $\displaystyle\int_0^{+\infty}\frac{dt}{(1+e^t)(1+e^{-t})}$. En multipliant numérateur et dénominateur par $e^t$ : $(1+e^{-t})e^t=e^t+1$, donc
\[
\frac{1}{(1+e^t)(1+e^{-t})}=\frac{e^t}{(1+e^t)^2},
\]
qui est de la forme $\dfrac{u'}{u^2}$ avec $u=1+e^t$ ; une primitive est $-\dfrac{1}{1+e^t}$. Ainsi
\[
\int_0^x\frac{e^t}{(1+e^t)^2}dt=\Big[-\frac1{1+e^t}\Big]_0^x=\frac12-\frac1{1+e^x}\xrightarrow[x\to+\infty]{}\frac12,\qquad\boxed{\int_0^{+\infty}\frac{dt}{(1+e^t)(1+e^{-t})}=\frac12}.
\]

\textbf{b.} $\displaystyle\int_1^{+\infty}\frac{e^{-\sqrt t}}{\sqrt t}dt$. Avec $u=-\sqrt t$, $u'=-\dfrac1{2\sqrt t}$, donc $\dfrac{e^{-\sqrt t}}{\sqrt t}=-2u'e^u$, de primitive $-2e^{-\sqrt t}$ :
\[
\int_1^x\frac{e^{-\sqrt t}}{\sqrt t}dt=\Big[-2e^{-\sqrt t}\Big]_1^x=\frac2e-2e^{-\sqrt x}\xrightarrow[x\to+\infty]{}\frac2e,\qquad\boxed{\int_1^{+\infty}\frac{e^{-\sqrt t}}{\sqrt t}dt=\frac2e}.
\]

\textbf{c.} $\displaystyle\int_0^{+\infty}\frac{\arctan t}{1+t^2}dt$. C'est de la forme $u'u$ avec $u=\arctan t$, de primitive $\frac12(\arctan t)^2$ :
\[
\int_0^x\frac{\arctan t}{1+t^2}dt=\frac{(\arctan x)^2}{2}\xrightarrow[x\to+\infty]{}\frac12\Big(\frac\pi2\Big)^2,\qquad\boxed{\int_0^{+\infty}\frac{\arctan t}{1+t^2}dt=\frac{\pi^2}{8}}.
\]

\textbf{d.} $\displaystyle\int_1^{+\infty}\frac{\ln t}{t^2}dt$. Par parties avec $u=\ln t$, $v'=\frac1{t^2}$ ($u'=\frac1t$, $v=-\frac1t$) :
\[
\int_1^x\frac{\ln t}{t^2}dt=\Big[-\frac{\ln t}{t}\Big]_1^x+\int_1^x\frac{dt}{t^2}=-\frac{\ln x}{x}+1-\frac1x\xrightarrow[x\to+\infty]{}1
\]
(croissances comparées : $\frac{\ln x}{x}\to0$). Donc $\boxed{\displaystyle\int_1^{+\infty}\frac{\ln t}{t^2}dt=1}$.

\textbf{e.} $\displaystyle\int_a^{+\infty}\frac{dt}{t(t+b)}$, $a>0$, $b>0$. On décompose : $\dfrac1{t(t+b)}=\dfrac1b\Big(\dfrac1t-\dfrac1{t+b}\Big)$. Pour $x>a$ :
\[
\int_a^x\frac{dt}{t(t+b)}=\frac1b\Big[\ln\frac{t}{t+b}\Big]_a^x=\frac1b\Big(\ln\frac{x}{x+b}-\ln\frac{a}{a+b}\Big)\xrightarrow[x\to+\infty]{}\frac1b\Big(0-\ln\frac{a}{a+b}\Big),
\]
car $\frac{x}{x+b}\to1$. Donc
\[
\boxed{\int_a^{+\infty}\frac{dt}{t(t+b)}=\frac1b\ln\frac{a+b}{a}=\frac1b\ln\Big(1+\frac ba\Big)}.
\]

\textbf{2. Nature des intégrales.}

\textbf{a)} $\displaystyle\int_0^{+\infty}\frac{\ln t}{t^2+1}dt$. La fonction est continue sur $]0,+\infty[$ : deux bornes à étudier.
\begin{itemize}
\item \emph{En $0$} : $\dfrac{\ln t}{t^2+1}\sim\ln t$, de signe constant (négatif) sur $]0,1]$. Or $\int_x^1\ln t\,dt=\big[t\ln t-t\big]_x^1=-1-x\ln x+x\to-1$ quand $x\to0^+$ : $\int_0^1\ln t\,dt$ converge. Par équivalence (fonctions de signe constant), $\int_0^1\frac{\ln t}{t^2+1}dt$ converge.
\item \emph{En $+\infty$} : $t^{3/2}\cdot\dfrac{\ln t}{t^2+1}\sim\dfrac{\ln t}{\sqrt t}\to0$, donc pour $t$ assez grand, $0\le\dfrac{\ln t}{t^2+1}\le\dfrac{1}{t^{3/2}}$. Comme $\int_1^{+\infty}\frac{dt}{t^{3/2}}$ converge (Riemann, $\frac32>1$), $\int_1^{+\infty}\frac{\ln t}{t^2+1}dt$ converge.
\end{itemize}
\textbf{L'intégrale converge.} (On peut montrer, par le changement de variable $t=\frac1u$, qu'elle vaut $0$.)

\textbf{b)} $\displaystyle\int_0^{+\infty}\frac{1-\cos t}{t^2}dt$. La fonction est continue et positive sur $]0,+\infty[$.
\begin{itemize}
\item \emph{En $0$} : $1-\cos t\sim\frac{t^2}{2}$, donc $\dfrac{1-\cos t}{t^2}\to\dfrac12$ : la fonction se prolonge par continuité en $0$, l'intégrale sur $]0,1]$ est « faussement impropre », donc convergente.
\item \emph{En $+\infty$} : $0\le\dfrac{1-\cos t}{t^2}\le\dfrac{2}{t^2}$ et $\int_1^{+\infty}\frac{2}{t^2}dt$ converge.
\end{itemize}
\textbf{L'intégrale converge} (elle vaut $\frac\pi2$).

\textbf{c)} $\displaystyle\int_0^{+\infty}\frac{\sin t+\cos t}{\sqrt{t^3+2}}dt$. La fonction est continue sur $[0,+\infty[$ ($t^3+2\ge2>0$) : seule la borne $+\infty$ pose problème. Pour $t\ge1$ :
\[
\Big|\frac{\sin t+\cos t}{\sqrt{t^3+2}}\Big|\le\frac{2}{\sqrt{t^3+2}}\le\frac{2}{t^{3/2}},
\]
et $\int_1^{+\infty}\frac{2}{t^{3/2}}dt$ converge. L'intégrale est \textbf{absolument convergente, donc convergente}. (La majoration ne doit être utilisée que sur $[1,+\infty[$ : $\int_0^1\frac{dt}{t^{3/2}}$ diverge.)

\textbf{d)} $\displaystyle\int_{-\infty}^{+\infty}\sin t\,dt$. Elle converge si et seulement si $\int_{-\infty}^0$ et $\int_0^{+\infty}$ convergent. Or $\int_0^x\sin t\,dt=1-\cos x$ n'a pas de limite quand $x\to+\infty$ (elle prend les valeurs $0$ et $2$ pour $x=2k\pi$ et $x=(2k+1)\pi$). Donc $\int_0^{+\infty}\sin t\,dt$ diverge et \textbf{$\int_{-\infty}^{+\infty}\sin t\,dt$ diverge}.

\remarque{On a pourtant $\int_{-x}^x\sin t\,dt=0$ pour tout $x$ (fonction impaire) : la « valeur principale » existe, mais cela ne suffit pas pour la convergence ; les deux bornes doivent être traitées indépendamment.}

\textbf{3. Discussion suivant $\alpha$.}

\textbf{a.} $\displaystyle\int_0^{+\infty}\frac{dt}{t^\alpha}$. Rappel (Riemann) : $\int_0^1\frac{dt}{t^\alpha}$ converge si et seulement si $\alpha<1$, et $\int_1^{+\infty}\frac{dt}{t^\alpha}$ converge si et seulement si $\alpha>1$. Ces deux conditions sont incompatibles : \textbf{$\int_0^{+\infty}\frac{dt}{t^\alpha}$ diverge pour tout $\alpha\in\mathbb{R}$}.

\textbf{b.} $\displaystyle\int_0^{+\infty}\frac{t-\sin t}{t^\alpha}dt$. Pour $t>0$, $t-\sin t>0$ : la fonction est continue et positive sur $]0,+\infty[$, et on peut utiliser les équivalents.
\begin{itemize}
\item \emph{En $+\infty$} : $t-\sin t\sim t$ (car $\frac{\sin t}{t}\to0$), donc $\dfrac{t-\sin t}{t^\alpha}\sim\dfrac1{t^{\alpha-1}}$ : converge si et seulement si $\alpha-1>1$, soit $\alpha>2$.
\item \emph{En $0$} : $\sin t=t-\frac{t^3}{6}+o(t^3)$, donc $t-\sin t\sim\frac{t^3}{6}$ et $\dfrac{t-\sin t}{t^\alpha}\sim\dfrac{1}{6\,t^{\alpha-3}}$ : converge si et seulement si $\alpha-3<1$, soit $\alpha<4$.
\end{itemize}
\[
\boxed{\int_0^{+\infty}\frac{t-\sin t}{t^\alpha}dt\ \text{converge}\iff2<\alpha<4}.
\]

\medskip\textbf{Exercice 2}

\textbf{1) Probabilité d'au moins une panne en 40 jours.}

Chaque jour, une machine donnée tombe en panne avec la probabilité $p=0{,}002$, indépendamment des autres jours (la machine est remise en service le jour ouvrable suivant, dans les mêmes conditions). L'évènement contraire de « au moins une panne » est « aucune panne pendant les 40 jours », c'est-à-dire « pas de panne le 1\textsuperscript{er} jour » et ... et « pas de panne le 40\textsuperscript{e} jour ». Par indépendance :
\[
P(\text{aucune panne})=(1-p)^{40}=0{,}998^{40}\approx0{,}92304,
\]
\[
P(\text{au moins une panne})=1-0{,}998^{40}\approx\boxed{0{,}077}.
\]

\textbf{2) Loi de $X$.} On répète $n=40$ fois, de façon indépendante et dans des conditions identiques, une épreuve de Bernoulli (une journée) dont le « succès » est « la machine tombe en panne », de probabilité $p=0{,}002$. $X$ compte le nombre de succès : \textbf{$X$ suit la loi binomiale $\mathcal{B}(40\,;\,0{,}002)$} :
\[
P(X=k)=C_{40}^k\,(0{,}002)^k\,(0{,}998)^{40-k},\qquad k\in\{0,\dots,40\}.
\]

\textbf{3) Même probabilité à l'aide de la loi de $X$.}
\[
P(X\ge1)=1-P(X=0)=1-C_{40}^0(0{,}002)^0(0{,}998)^{40}\approx1-0{,}92304=\boxed{0{,}077}.
\]
On retrouve le résultat de la question 1.

\remarque{Comme $n$ est grand et $p$ petit ($np=0{,}08$), on peut approcher $X$ par une loi de Poisson de paramètre $\lambda=np=0{,}08$ : $P(X\ge1)\approx1-e^{-0{,}08}\approx0{,}0769$. Par ailleurs, $P(X=1)=40\times0{,}002\times0{,}998^{39}\approx0{,}07399$, d'où $P(X\ge2)\approx0{,}00297$ : les pannes multiples sont très rares.}

\textbf{4) Loi, espérance, variance et écart-type de $Y$.}

Les 52 machines fonctionnent de façon indépendante et dans des conditions identiques ; chacune tombe au moins une fois en panne en 40 jours avec la probabilité $p'=0{,}077$. $Y$ compte le nombre de machines concernées : $Y$ suit la loi binomiale $\mathcal{B}(52\,;\,0{,}077)$. Donc
\[
E(Y)=np'=52\times0{,}077=4{,}004\approx4,\qquad V(Y)=np'(1-p')=4{,}004\times0{,}923\approx3{,}696,
\]
\[
\sigma(Y)=\sqrt{V(Y)}\approx1{,}92 .
\]
En moyenne, environ 4 photocopieurs du parc tombent en panne au moins une fois par période de 40 jours.

\medskip\textbf{Exercice 3}

$f(x,y)=\dfrac{2x^3+3xy^2}{x^2+y^2}$ si $(x,y)\neq(0,0)$, $f(0,0)=0$.

\textbf{1. Continuité.} Sur $\mathbb{R}^2\setminus\{(0,0)\}$, $f$ est une fraction rationnelle dont le dénominateur ne s'annule pas : elle y est continue.

En $(0,0)$ : en coordonnées polaires $x=r\cos\theta$, $y=r\sin\theta$ ($r>0$),
\[
f(x,y)=\frac{r^3(2\cos^3\theta+3\cos\theta\sin^2\theta)}{r^2}=r\cos\theta\,(2\cos^2\theta+3\sin^2\theta).
\]
Donc $|f(x,y)|\le r(2+3)=5r=5\sqrt{x^2+y^2}$, majoration \emph{indépendante de $\theta$}, qui tend vers $0$ quand $(x,y)\to(0,0)$. Ainsi $\lim_{(0,0)}f=0=f(0,0)$ : $f$ est continue en $(0,0)$, donc sur $\mathbb{R}^2$.

\textbf{2. Dérivées directionnelles en $(0,0)$ et gradient.}

Soit $u=(u_1,u_2)$ un vecteur unitaire ($u_1^2+u_2^2=1$). Pour $t\neq0$ :
\[
\frac{f(tu_1,tu_2)-f(0,0)}{t}=\frac1t\cdot\frac{t^3(2u_1^3+3u_1u_2^2)}{t^2(u_1^2+u_2^2)}=2u_1^3+3u_1u_2^2 .
\]
Ce quotient est constant, donc a une limite quand $t\to0$ : $f$ admet en $(0,0)$ une dérivée dans toutes les directions, et
\[
\boxed{D_uf(0,0)=2u_1^3+3u_1u_2^2=u_1(2+u_2^2)}.
\]
En particulier (directions $(1,0)$ et $(0,1)$) : $\dfrac{\partial f}{\partial x}(0,0)=2$ et $\dfrac{\partial f}{\partial y}(0,0)=0$, d'où
\[
\boxed{\overrightarrow{\operatorname{grad}}f(0,0)=(2,0)}.
\]

\textbf{3. Différentiabilité en $(0,0)$.}

Si $f$ était différentiable en $(0,0)$, sa différentielle serait $(h,k)\mapsto2h$ (donnée par le gradient), et l'on aurait $D_uf(0,0)=\overrightarrow{\operatorname{grad}}f(0,0)\cdot u=2u_1$ pour tout $u$. Or, pour $u=\big(\frac{1}{\sqrt2},\frac1{\sqrt2}\big)$ : $D_uf(0,0)=\frac1{\sqrt2}\big(2+\frac12\big)=\frac{5}{2\sqrt2}\neq\frac{2}{\sqrt2}$.

Autrement dit, $f(x,y)-f(0,0)-2x=\dfrac{2x^3+3xy^2-2x^3-2xy^2}{x^2+y^2}=\dfrac{xy^2}{x^2+y^2}=r\cos\theta\sin^2\theta$, qui n'est pas un $o(r)$ (pour $\theta=\frac\pi4$, il vaut $\frac{r}{2\sqrt2}$).

\textbf{$f$ n'est pas différentiable en $(0,0)$}, bien qu'elle y soit continue et y admette des dérivées dans toutes les directions.

\medskip\textbf{Exercice 4}

$f(x)=\displaystyle\int_0^1\frac{t^{x-1}}{1+t}dt$.

\textbf{1. Domaine de définition.} Pour $x$ fixé, $t\mapsto\frac{t^{x-1}}{1+t}$ est continue et positive sur $]0,1]$ ; le seul problème est en $0$, où $\dfrac{t^{x-1}}{1+t}\sim t^{x-1}=\dfrac{1}{t^{1-x}}$. D'après le critère de Riemann et le critère d'équivalence (fonctions positives), l'intégrale converge si et seulement si $1-x<1$, c'est-à-dire $x>0$.
\[
\boxed{\mathcal{D}_f=]0,+\infty[}.
\]

\textbf{2. Continuité.} On applique le théorème de continuité des intégrales à paramètre sur tout segment $[a,b]\subset]0,+\infty[$ :
\begin{itemize}
\item pour tout $t\in]0,1]$, $x\mapsto\frac{t^{x-1}}{1+t}=\frac{e^{(x-1)\ln t}}{1+t}$ est continue ;
\item pour tout $x$, $t\mapsto\frac{t^{x-1}}{1+t}$ est continue par morceaux sur $]0,1]$ ;
\item \emph{domination} : pour $x\in[a,b]$ et $t\in]0,1]$, comme $\ln t\le0$ et $x-1\ge a-1$, on a $(x-1)\ln t\le(a-1)\ln t$, donc $0\le\dfrac{t^{x-1}}{1+t}\le t^{a-1}=\varphi(t)$, et $\int_0^1\varphi$ converge car $a>0$.
\end{itemize}
Donc $f$ est continue sur $[a,b]$ pour tous $0<a<b$, et donc sur $]0,+\infty[$.

\textbf{3. $f(x)+f(x+1)$ pour $x>0$.}
\[
f(x)+f(x+1)=\int_0^1\frac{t^{x-1}+t^x}{1+t}dt=\int_0^1\frac{t^{x-1}(1+t)}{1+t}dt=\int_0^1t^{x-1}dt=\Big[\frac{t^x}{x}\Big]_0^1=\boxed{\frac1x}.
\]

\textbf{4. Équivalent en $0$.} Pour $x>0$, $f(x)=\frac1x-f(x+1)$, donc $xf(x)=1-xf(x+1)$. Comme $f$ est continue en $1$, $f(x+1)\to f(1)=\int_0^1\frac{dt}{1+t}=\ln2$ quand $x\to0^+$, et donc $xf(x+1)\to0$. Ainsi $xf(x)\to1$ :
\[
\boxed{f(x)\underset{x\to0^+}{\sim}\frac1x}.
\]

\textbf{5. Limite en $+\infty$.} Pour $t\in]0,1]$, $\frac{1}{1+t}\le1$, donc
\[
0\le f(x)\le\int_0^1t^{x-1}dt=\frac1x\xrightarrow[x\to+\infty]{}0 .
\]
Par encadrement, $\boxed{\lim\limits_{x\to+\infty}f(x)=0}$. (On peut aussi en déduire $f(x)\sim\frac{1}{2x}$ en $+\infty$, car $f$ est décroissante et $f(x)+f(x+1)=\frac1x$.)

\end{spacing}
\end{document}
